\chapter{INTRODUCTION}

\begingroup
\hyphenpenalty=10000
\exhyphenpenalty=10000
\sloppy

The rapid expansion of cloud-native applications and microservices-based architectures has significantly increased the complexity of resource management in distributed systems \cite{khaleq2021, fourati2021}. Modern applications, particularly those handling highly dynamic and bursty workloads such as university course registration platforms, experience sudden and unpredictable traffic spikes. These fluctuations can result in severe performance degradation, increased response times, and inefficient resource utilization if not managed effectively. Consequently, ensuring consistent performance under varying workloads has become a critical challenge in cloud environments.

Traditional autoscaling mechanisms primarily rely on static CPU and memory utilization thresholds to trigger scaling decisions \cite{khaleq2021}. Although simple to implement, these approaches often fail to capture actual system behavior under load, particularly in scenarios involving request queuing and service dependencies. As a result, they tend to react too late to traffic surges or scale inefficiently, leading to over-provisioning or under-provisioning of resources. Moreover, these methods treat system metrics independently, ignoring the underlying relationships between workload arrival rates, service capacity, and response time.

Recent advancements in performance modeling and queuing theory offer a more accurate and analytical framework for understanding system behavior \cite{wang2024}. By modeling services as queuing systems, it becomes possible to estimate critical performance metrics such as utilization, waiting time, and overall response time. However, most existing autoscaling solutions do not fully leverage these theoretical models in real-time decision-making, thereby limiting their ability to perform proactive and optimal scaling.

To address these limitations, this project presents an intelligent, queuing-theory-driven autoscaling framework for a microservices-based course registration system. The proposed approach models each service as an M/M/c queue and continuously analyzes real-time metrics, including request arrival rates and service latency. A Multi-Feature Decision System (MFDS) evaluates system performance by estimating response time and comparing it against predefined Service Level Agreement (SLA) thresholds to determine appropriate scaling actions \cite{lai2025}.

To further improve efficiency, a Minimum Scaling Overhead (MSO) mechanism computes the optimal number of service replicas directly, thereby avoiding incremental scaling delays. Additionally, a Flexible Probability Routing (FPR) strategy dynamically distributes incoming traffic across services based on both workload demand and available capacity, ensuring balanced load distribution \cite{arulmary2026}. The system is implemented using a containerized microservices architecture orchestrated via Docker Swarm and monitored using Prometheus, enabling real-time scaling decisions and performance optimization under dynamic workloads.

\section{BACKGROUND}

Modern cloud-native applications and enterprise systems expose multiple microservices to users, applications, and external clients, significantly increasing the complexity of traffic management and resource allocation \cite{fourati2021}. In high-demand scenarios, such as course registration systems, sudden bursts of incoming requests can lead to congestion, increased queuing delays, and uneven load distribution across services. These workload patterns often include rapid spikes in request arrival rates, short-lived bursts of traffic, and fluctuating service demand, making efficient resource management a challenging task.

Conventional autoscaling solutions struggle to handle such dynamic behaviors effectively \cite{khaleq2021}. Threshold-based approaches, which rely on CPU or memory utilization, fail to capture the actual performance state of the system, particularly the queuing effects that directly impact response time. These methods often react too late to workload changes or scale inefficiently, leading to increased latency during peak demand, and resource wastage during low utilization periods. Additionally, traditional approaches treat services independently, ignoring the interdependencies and traffic flow relationships within microservice architectures.

Queuing theory provides a more accurate framework for modeling system behavior by representing each service as a queuing system, capturing the relationship between request arrival rates, service rates, and system utilization \cite{lai2025}. By analyzing these parameters, it becomes possible to estimate key performance metrics such as waiting time, queue length, and overall response time.

To overcome these limitations, the proposed system leverages a queuing-theory-driven approach combined with intelligent decision mechanisms. Each microservice is modeled as an M/M/c queue, and real-time metrics are continuously analyzed to evaluate system performance. A Multi-Feature Decision System (MFDS) determines scaling actions based on predicted response time and SLA constraints. This is further enhanced by a Minimum Scaling Overhead (MSO) strategy for optimal replica allocation and a Flexible Probability Routing (FPR) mechanism to ensure efficient and balanced traffic distribution across services, resulting in improved system stability and performance under dynamic workloads.


\section{PROBLEM STATEMENT}

Traditional autoscaling methods based on CPU and memory thresholds fail to handle dynamic workloads effectively, leading to delayed scaling and inefficient resource utilization \cite{khaleq2021}. These reactive approaches are particularly inadequate for bursty traffic patterns commonly observed in university course registration systems, where sudden enrollment periods can generate traffic spikes that exceed normal capacity by several orders of magnitude within minutes.

The limitations of threshold-based autoscaling become more pronounced in microservices architectures, where service dependencies create cascading effects that amplify performance degradation. When one service becomes overwhelmed, it creates bottlenecks that propagate throughout the entire system, resulting in increased latency, timeout errors, and poor user experience during critical registration periods.

The time required to detect threshold violations, initiate scaling actions, and provision new instances often exceeds the duration of traffic spikes, rendering the scaling response ineffective. This reactive nature also leads to resource over-provisioning during scale-down phases, as conservative policies are employed to avoid service disruption.

There is a need for an intelligent autoscaling system that leverages queuing theory and real-time metrics to make accurate, proactive scaling decisions while maintaining SLA requirements. Such a system must be capable of predicting resource demands based on current system state, optimizing scaling decisions to minimize both response time and resource overhead, and dynamically distributing load across available service instances to ensure balanced utilization and improved overall system performance.

\section{CHALLENGES}

Managing dynamic workloads in microservices-based systems presents several challenges in real-world cloud environments. Sudden traffic spikes and fluctuating request patterns make it difficult for traditional autoscaling methods to respond effectively, leading to increased latency and potential service degradation. This problem is particularly acute in university course registration systems, where enrollment periods create highly concentrated traffic bursts that can overwhelm system capacity within minutes, resulting in student frustration and administrative complications.

Conventional autoscaling approaches rely on static CPU and memory thresholds, which do not accurately reflect system performance or queuing behavior \cite{fourati2021}. These metrics provide limited insight into the actual service capacity and fail to capture the complex relationship between request arrival rates, processing times, and queue lengths. This results in delayed or inefficient scaling decisions, especially under bursty workloads where the system may experience significant performance degradation before threshold violations are detected and scaling actions are initiated.

The reactive nature of threshold-based systems introduces additional latency through the scaling process itself. By the time new instances are provisioned and become available, traffic spikes may have already subsided, leading to unnecessary resource allocation and increased operational costs. Conversely, during sustained high-load periods, the incremental scaling approach often fails to provision sufficient resources quickly enough, resulting in prolonged service degradation.

Additionally, microservices architectures involve interdependent services, and treating them independently can lead to uneven load distribution and resource imbalance. In a course registration system, the authentication service, course catalog service, and seat reservation service form a processing chain where bottlenecks in any component can cascade throughout the entire system. Traditional autoscaling methods lack the holistic view necessary to coordinate scaling decisions across these interconnected services, often resulting in scenarios where some services are over-provisioned while others remain under-resourced.

Scalability is another challenge, as the system must efficiently handle high volumes of requests while ensuring optimal resource utilization and maintaining Service Level Agreements (SLAs). The system must balance competing objectives: minimizing response times to ensure user satisfaction, optimizing resource costs to maintain operational efficiency, and preserving system stability during rapid scaling operations. Furthermore, the heterogeneous nature of microservices, each with different resource requirements and performance characteristics, complicates the development of unified scaling strategies that can effectively manage the entire system ecosystem.

\section{MOTIVATION}

The increasing demand for highly scalable and responsive cloud-based applications highlights the need for more intelligent autoscaling mechanisms \cite{wang2024}. Modern applications, particularly those serving critical functions like university course registration systems, must handle unpredictable traffic patterns that can vary by orders of magnitude within short time periods. Traditional threshold-based approaches are insufficient for handling dynamic and bursty workloads, often resulting in performance degradation and inefficient resource utilization. These conventional methods fail to capture the complex interdependencies between system components and lack the predictive capabilities necessary to anticipate resource demands before performance issues manifest.

The limitations of reactive autoscaling become particularly evident during peak usage periods, such as course enrollment deadlines, where thousands of students attempt to register simultaneously. In such scenarios, the delay inherent in threshold-based detection and response mechanisms can lead to system overload, resulting in timeouts, failed transactions, and poor user experience. Moreover, the incremental nature of traditional scaling often proves inadequate when faced with sudden traffic surges, as the time required to provision additional resources exceeds the duration needed to handle the workload effectively.

Accurate modeling of system behavior using queuing theory provides a more reliable way to predict performance metrics such as response time and utilization \cite{lai2025}. By representing each microservice as a queuing system with specific arrival rates, service rates, and server configurations, it becomes possible to mathematically predict system behavior under various load conditions. This theoretical foundation enables the development of proactive scaling algorithms that can anticipate resource requirements based on current system state and incoming traffic patterns, rather than merely reacting to performance degradation after it occurs.

Integrating these models with real-time decision-making enables proactive and precise scaling actions, improving overall system stability. The combination of queuing theory with modern monitoring capabilities allows for the development of sophisticated decision systems that can evaluate multiple performance indicators simultaneously, considering factors such as current queue lengths, service utilization, and predicted arrival rates. This multi-dimensional approach enables more nuanced scaling decisions that optimize not only for immediate performance requirements but also for long-term resource efficiency and cost effectiveness.

Furthermore, the distributed nature of microservices architectures necessitates coordination mechanisms that can manage scaling decisions across multiple interdependent services. A holistic approach that considers the entire service chain, from authentication through course selection to seat reservation, ensures that scaling actions are coordinated to prevent bottlenecks from shifting between services rather than being eliminated.

This project is motivated by the objective of designing a scalable, efficient, and intelligent autoscaling framework that bridges the gap between theoretical queuing models and practical cloud deployment systems, ensuring optimal resource utilization while maintaining Service Level Agreement (SLA) requirements. The framework aims to demonstrate that mathematical modeling can be effectively translated into operational autoscaling systems that outperform traditional approaches in terms of response time, resource efficiency, and user satisfaction. By implementing this system in a realistic course registration scenario, the project seeks to validate the practical applicability of queuing theory-based autoscaling in production cloud environments.

\section{PROPOSED WORK}

This project proposes a queuing-theory-based intelligent autoscaling system for microservices that addresses the limitations of traditional threshold-based approaches through mathematical modeling and predictive analytics. Each service is modeled as an M/M/c queue, representing a multi-server queuing system where requests arrive according to a Poisson process and are served with exponentially distributed service times. This mathematical foundation enables accurate prediction of key performance metrics including system utilization, average queue length, and expected response time under varying load conditions \cite{lai2025}.

Real-time metrics collected from the microservices environment are used to make scaling decisions using a Multistage Fine-Grained Dynamic Scaling (MFDS) mechanism that evaluates multiple performance indicators simultaneously. The MFDS approach considers not only current system state but also predicted future demands based on traffic patterns and historical data, enabling proactive scaling decisions that prevent performance degradation before it occurs \cite{khaleq2021}.

A Microservice Scaling Optimization (MSO) strategy computes the optimal number of service replicas required to meet Service Level Agreement (SLA) requirements while minimizing resource overhead. Unlike traditional incremental scaling approaches that add or remove instances one at a time, the MSO algorithm calculates the target replica count in a single step, significantly reducing scaling latency and improving system responsiveness during traffic spikes \cite{wang2024}.

To complement the scaling mechanisms, Fair Probabilistic Routing (FPR) ensures efficient and balanced load distribution across available service instances. The FPR algorithm dynamically adjusts routing probabilities based on both current service capacity and real-time performance metrics, preventing hotspots and ensuring optimal utilization of all available resources \cite{arulmary2026}. This approach is particularly effective in microservices architectures where service interdependencies can create cascading performance issues if load distribution is not carefully managed.

The system is implemented using Docker Swarm for container orchestration and Prometheus for real-time monitoring and metrics collection, creating a comprehensive autoscaling framework that integrates seamlessly with modern cloud-native deployment pipelines. The containerized approach provides the flexibility and scalability necessary for dynamic resource management while maintaining system stability during scaling operations \cite{fourati2021}.

Experimental evaluation demonstrates that the proposed approach achieves improved performance and resource utilization under dynamic workloads, with significant reductions in response time and scaling latency compared to conventional threshold-based autoscaling methods. The integration of queuing theory with real-time decision-making enables more accurate predictions of system behavior, resulting in proactive scaling actions that maintain optimal performance while minimizing resource waste \cite{lai2025}.


\section{ORGANIZATION OF THE REPORT}

The report is systematically organized into chapters to provide a clear and structured presentation of the project. Each chapter focuses on a specific aspect of the work, covering the background, system design, implementation, and evaluation. The organization of the report is as follows:

\vspace{0.5cm}
\textbf{Chapter 1} presents an introduction to intelligent autoscaling in microservices-based systems, establishing the context for cloud-native applications and their scaling challenges. This chapter includes the problem statement addressing limitations of traditional threshold-based approaches, identifies key challenges in dynamic workload management, discusses the motivation for queuing theory-based solutions, outlines the proposed work with its core contributions, and provides the overall report structure.
\\

\vspace{0.4cm}
\textbf{Chapter 2} provides a comprehensive review of existing autoscaling techniques and related work in the field. This includes an analysis of threshold-based methods and their reactive limitations, examination of cloud resource management strategies and container orchestration approaches, exploration of the application of queuing theory in performance modeling and prediction, review of machine learning approaches to autoscaling, and identification of research gaps that justify the need for the proposed intelligent autoscaling framework.

\vspace{0.4cm}
\textbf{Chapter 3} describes the overall architecture and design methodology of the proposed autoscaling system. This chapter covers the microservices architecture design for the course registration platform, detailed explanation of the queuing model (M/M/c) and its mathematical foundations, comprehensive description of the core components including Multistage Fine-Grained Dynamic Scaling (MFDS) for decision-making, Microservice Scaling Optimization (MSO) for replica calculation, and Fair Probabilistic Routing (FPR) for load distribution, along with the integration strategy for real-time monitoring and control loops.

\vspace{0.4cm}
\textbf{Chapter 4} discusses the comprehensive implementation details of the system, covering the development of individual microservices (authentication, course catalog, and seat reservation services), integration with Docker Swarm for container orchestration and dynamic scaling, implementation of Prometheus-based monitoring infrastructure for real-time metrics collection, detailed autoscaler service design with queuing theory algorithms, configuration of the experimental environment, and development of traffic simulation tools for performance evaluation.\\

\vspace{0.4cm}
\textbf{Chapter 5} presents the evaluation and performance analysis of the proposed system through extensive experimentation. This includes detailed response time analysis under various load conditions, examination of scaling behavior and decision accuracy, assessment of resource utilization efficiency and cost optimization, comparative analysis with traditional threshold-based autoscaling approaches, evaluation of system stability and SLA compliance, and discussion of the practical implications and benefits of the queuing theory-based approach.

\vspace{0.4cm}
\textbf{Chapter 6} concludes the report by summarizing the key findings and contributions of the research, discussing the practical significance of the results for cloud-native application deployment, identifying limitations of the current approach, and outlining potential future enhancements including advanced machine learning integration, multi-cloud deployment strategies, and improvements in intelligent autoscaling systems for next-generation microservices architectures.

\endgroup